\uuid{mPSm}
\titre{Éléments simples}
\auteur{bodin ; hanani ; mzari}
\organisation{exo7}
\licence{CC BY-NC-SA 4.0}
\datecreate{2020-01-20}
\niveau{L1}
\module{Algèbre}
\chapitre{Polynôme, fraction rationnelle}
\sousChapitre{Fraction rationnelle}
\difficulte{1}
\theme{Fractions rationnelles, éléments simples}
\exo7qcmid{lille-1.polynomes.026}

\contenu{
\begin{question}
\qclassification{105.04}
\qtags{motcle=éléments simples}
Soient $P(X)=X-1$, $Q(X)=(X+1)^2(X^2+X+1)$. On décompose la fraction $F = \frac{P}{Q}$ sur $\Rr$.
\begin{answers}
    \good{La partie polynomiale est nulle.}

    \bad{Il peut y avoir un élément simple $\frac{a}{X-1}$.}

    \bad{Il peut y avoir un élément simple $\frac{a}{X+1}$ mais pas  $\frac{a}{(X+1)^2}$.}

    \good{Il peut y avoir un élément simple $\frac{aX+b}{X^2+X+1}$ mais pas  $\frac{aX+b}{(X^2+X+1)^2}$.}

\end{answers}
\begin{explanations}
$\frac{P(X)}{Q(X)} = \frac{X-1}{(X+1)^2(X^2+X+1)}
= \frac{-1}{X+1}+\frac{-2}{(X+1)^2}+\frac{X+2}{X^2+X+1}$.
\end{explanations}
\end{question}
}
